\subsection{Lösung und Lösungsmenge}
\label{subsec:linear_systems_solution_sets_and_geometry}

\begin{DefinitionBox}{Lösung und Lösungsmenge}{linear_systems_solution_set}
  Ein Tupel $(x_1,\dots,x_n)\in\R^n$ heißt \emph{Lösung} eines LGS, wenn es
  alle Gleichungen gleichzeitig erfüllt. Die Menge aller Lösungen heißt
  \emph{Lösungsmenge} $L$.
\end{DefinitionBox}

Ein lineares Gleichungssystem besitzt entweder genau eine, keine oder
unendlich viele Lösungen. Für zwei Unbekannte lässt sich dies geometrisch
durch Geraden veranschaulichen:

\begin{center}
  \begin{tikzpicture}[
  >=Latex,
  axis/.style={->,black!60,line width=0.5pt},
  linea/.style={THMBlue,line width=1.1pt},
  lineb/.style={black!60,line width=1.1pt},
  title/.style={font=\small\bfseries,align=center},
  result/.style={font=\small,align=center}
]
  \foreach \shift/\case in {0/one,5.5/none,11/many}{
    \begin{scope}[xshift=\shift cm]
      \draw[axis] (-2,-1.65)--(2,-1.65) node[right] {$x_1$};
      \draw[axis] (0,-2)--(0,1.75) node[above] {$x_2$};
    \end{scope}
  }

  \begin{scope}
    \draw[linea] (-1.75,-1.95)--(1.55,1.35);
    \draw[lineb] (-1.65,1.25)--(1.75,-2.15);
    \fill[THMBlue] (0,-0.2) circle (1.8pt);
    \node[title] at (0,2.35) {genau eine Lösung};
    \node[result] at (0,-2.65) {$L=\{S\}$};
  \end{scope}

  \begin{scope}[xshift=5.5cm]
    \draw[linea] (-1.75,-1.75)--(1.7,1.7);
    \draw[lineb] (-1.75,-0.95)--(0.9,1.7);
    \node[title] at (0,2.35) {keine Lösung};
    \node[result] at (0,-2.65) {$L=\varnothing$};
  \end{scope}

  \begin{scope}[xshift=11cm]
    \draw[lineb,line width=2.6pt] (-1.75,-1.75)--(1.7,1.7);
    \draw[linea] (-1.75,-1.75)--(1.7,1.7);
    \node[title] at (0,2.35) {unendlich viele Lösungen};
    \node[result] at (0,-2.65) {$L=$ gemeinsame Gerade};
  \end{scope}
\end{tikzpicture}

\end{center}

Die Bilder zeigen nur den zweidimensionalen Fall. In höheren Dimensionen
beschreibt jede lineare Gleichung eine Hyperebene; die drei möglichen Anzahlen
von Lösungen bleiben dieselben.

\begin{KeyPointBox}{Nur drei Lösungstypen}{linear_systems_three_solution_types}
  Ein LGS über $\R$ besitzt
  \[
    \boxed{\text{keine}},\qquad
    \boxed{\text{genau eine}}
    \qquad\text{oder}\qquad
    \boxed{\text{unendlich viele Lösungen}}.
  \]
  Zwei verschiedene Lösungen erzeugen bereits eine ganze Gerade von Lösungen;
  daher kann es nicht genau zwei oder allgemein endlich viele Lösungen größer
  als eins geben.
\end{KeyPointBox}
